\documentclass[9pt,twoside]{extarticle}
\usepackage{geometry}
\usepackage{setspace}
\usepackage{fancyhdr}
\usepackage{textcase}
\usepackage{parskip}
\usepackage{fontspec}
\usepackage{microtype}
\usepackage{enumitem}

\usepackage{tikz}
\usetikzlibrary{calc, positioning, decorations.pathreplacing}

\geometry{
  paperwidth=6in,
  paperheight=9in,
  margin=0.75in,
}

\setstretch{1.15}
\setlength{\parindent}{1.25em}
\setlength{\parskip}{0pt}

\newcommand{\essaytitle}{The Essence of Tragedy}
\newcommand{\essaysubtitle}{and Other Footnotes and Papers}
\newcommand{\essayauthor}{Maxwell Anderson}

\title{\essaytitle}
\author{\essayauthor}

\makeatletter
\let\thetitle\@title
\let\theauthor\@author
\makeatother

\pagestyle{fancy}
\fancyhf{}
\fancyhead[LO,RE]{\thepage}
\fancyhead[CO]{\scriptsize\MakeTextUppercase{\essayauthor}}
\fancyhead[CE]{\scriptsize\MakeTextUppercase{\essaytitle}}
\renewcommand{\headrulewidth}{0pt}

\makeatletter
\renewcommand\section{\@startsection{section}{1}{\z@}%
  {-2.5ex \@plus -1ex \@minus -.2ex}%
  {1ex \@plus .2ex}%
  {\itshape\normalsize}}
\makeatother

\makeatletter
\renewcommand{\maketitle}{
  \thispagestyle{empty}
  \begin{center}
    {\fontseries{bx}\Large\MakeTextUppercase\thetitle\par}
    {\fontseries{bx}\large\essaysubtitle\par}
    {\footnotesize{\MakeTextUppercase{\theauthor}}}
  \end{center}
}
\makeatother

\begin{document}

\maketitle

\vspace{2em}


\vspace*{\fill}
\begin{quote}
THE ESSENCE OF TRAGEDY was written to be read at a session of The Modern Language Association in New York City, January, 1938. WHATEVER HOPE WE HAVE was delivered, in a slightly different form, as the Founder's Day Address at Carnegie Institute, 1937. A PRELUDE TO POETRY IN THE THEATRE and THE POLITICS OF KNICKERBOCKER HOLIDAY appeared as prefaces to published plays. YES, BY THE ETERNAL was printed in Stage as a rejoinder to an article by Max Eastman in the same magazine called ``By the Eternal."
\end{quote}
\vspace*{\fill}

\vspace{3em}
\tableofcontents

\vspace*{\fill}

\newpage


\section{The Essence of Tragedy}

Anybody who dares to discuss the making of tragedy lays himself open to critical assault and general barrage, for the theorists have been hunting for the essence of tragedy since Aristotle without entire success. There is no doubt that playwrights have occasionally written tragedy successfully, from Aeschylus on, and there is no doubt that Aristotle came very close to a definition of what tragedy is in his famous passage on catharsis. But why the performance of a tragedy should have a cleansing effect on the audience, why an audience is willing to listen to tragedy, why tragedy has a place in the education of men, has never, to my knowledge, been convincingly stated. I must begin by saying that I have not solved the Sphinx's riddle which fifty generations of skillful brains have left in shadow. But I have one suggestion which I think might lead to a solution if it were put to laboratory tests by those who know something about philosophical analysis and dialectic.

There seems no way to get at this suggestion except through a reference to my own adventures in playwriting, so I ask your tolerance while I use myself as an instance. A man who has written successful plays is usually supposed to know something about the theory of playwriting, and perhaps he usually does. In my own case, however, I must confess that I came into the theatre unexpectedly, without preparation, and stayed in it because I had a certain amount of rather accidental success. It was not until after I had fumbled my way through a good many successes and an appalling number of failures that I began to doubt the sufficiency of dramatic instinct and to wonder whether or not there were general laws governing dramatic structure which so poor a head for theory as my own might grasp and use. I had read the Poetics long before I tried playwriting, and I had looked doubtfully into a few well-known handbooks on dramatic structure, but the maxims and theories propounded always drifted by me in a luminous haze--brilliant, true, profound in context, yet quite without meaning for me when I considered the plan for a play or tried to clarify an emotion in dialogue. So far as I could make out every play was a new problem, and the old rules were inapplicable. There were so many rules, so many landmarks, so many pitfalls, so many essential reckonings, that it seemed impossible to find your way through the jungle except by plunging ahead, trusting to your sense of direction and keeping your wits about you as you went.

But as the seasons went by and my failures fell as regularly as the leaves in autumn I began to search again among the theorists of the past for a word of wisdom that might take some of the gamble out of playwriting. What I needed most of all, I felt, was a working definition of what a play is, or perhaps a formula which would include all the elements necessary to a play structure. A play is almost always, probably, an attempt to recapture a vision for the stage. But when you are working in the theatre it's most unsatisfactory to follow the gleam without a compass, quite risky to trust ``the light that never was on sea or land'' without making sure beforehand that you are not being led straight into a slough of despond. In other words you must make a choice among visions, and you must check your chosen vision carefully before assuming that it will make a play. But by what rules, what maps, what fields of reference can you check so intangible a substance as a revelation, a dream, an inspiration, or any similar nudge from the subconscious mind?

I shan't trouble you with the details of my search for a criterion, partly because I can't remember it in detail. But I re-read Aristotle's Poetics in the light of some bitter experience, and one of his observations led me to a comparison of ancient and modern playwriting methods. In discussing construction he made a point of the recognition scene as essential to tragedy. The recognition scene, as Aristotle isolated it in the tragedies of the Greeks, was generally an artificial device, a central scene in which the leading character saw through a disguise, recognized as a friend or as an enemy, perhaps as a lover or a member of his own family, some person whose identity had been hidden. Iphigeneia, for example, acting as priestess in an alien country, receives a victim for sacrifice and then recognizes her own brother in this victim. There is an instant and profound emotional reaction, instantly her direction in the play is altered. But occasionally, in the greatest of the plays, the recognition turned on a situation far more convincing, though no less contrived. Oedipus, hunting savagely for the criminal who has brought the plague upon Thebes, discovers that he is himself that criminal-- and since this is a discovery that affects not only the physical well-being and happiness of the hero, but the whole structure of his life, the effect on him and on the direction of the story is incalculably greater than could result from the more superficial revelation made to Iphigeneia.

Now scenes of exactly this sort are rare in the modern drama except in detective stories adapted for the stage. But when I probed a little more deeply into the memorable pieces of Shakespeare's theatre and our own I began to see that though modern recognition scenes are subtler and harder to find, they are none the less present in the plays we choose to remember. They seldom have to do with anything so naive as disguise or the unveiling of a personal identity. But the element of discovery is just as important as ever. For the mainspring in the mechanism of a modern play is almost invariably a discovery by the hero of some element in his environment or in his own soul of which he has not been aware--or which he has not taken sufficiently into account. Moreover, nearly every teacher of playwriting has had some inkling of this, though it was not until after I had worked out my own theory that what they said on this point took on accurate meaning for me. I still think that the rule which I formulated for my own guidance is more concise than any other, and so I give it here: A play should lead up to and away from a central crisis, and this crisis should consist in a discovery by the leading character which has an indelible effect on his thought and emotion and completely alters his course of action. The leading character, let me say again, must make the discovery; it must affect him emotionally; and it must alter his direction in the play.

Try that formula on any play you think worthy of study, and you will find that, with few exceptions, it follows this pattern or some variation of this pattern. The turning point of The Green Pastures, for example, is the discovery of God, who is the leading character, that even he must learn and grow, that a God who is to endure must conform to the laws of change. The turning point of Hamlet is Hamlet's discovery, in the play-scene, that his uncle was unquestionably the murderer of his father. In Abe Lincoln in Illinois Lincoln's discovery is that he has been a coward, that he has stayed out of the fight for the Union because he was afraid. In each case, you will note, the discovery has a profound emotional effect on the hero, and gives an entirely new direction to his action in the play. I'm not writing a disquisition on playwriting and wouldn't be competent to write one, but I do want to make a point of the superlative usefulness of this one touchstone for play-structure. When a man sets out to write a play his first problem is his subject and the possibilities of that subject as a story to be projected from the stage. His choice of subject matter is his personal problem, and one that takes its answer from his personal relation to his times. But if he wants to know a possible play subject when he finds it, if he wants to know how to mould the subject into play form after he has found it, I doubt that he'll ever discover another standard as satisfactory as the modern version of Aristotle which I have suggested. If the plot he has in mind does not contain a playable episode in which the hero or heroine makes an emotional discovery, a discovery that practically dictates the end of the story, then such an episode must be inserted--and if no place can be found for it the subject is almost certainly a poor one for the theatre. If this emotional discovery is contained in the story, but is not central, then it must be made central, and the whole action must revolve around it. In a three-act play it should fall near the end of the second act, though it may be delayed till the last; in a five-act play it will usually be found near the end of the third, though here also it can be delayed. Everything else in the play should be subordinated to this one episode--should lead up to or away from it.

Now this prime rule has a corollary which is just as important as the rule itself. The hero who is to make the central discovery in a play must not be a perfect man. He must have some variation of what Aristotle calls a tragic fault--and the reason he must have it is that when he makes his discovery he must change both in himself and in his action--and he must change for the better. The fault can be a very simple one--a mere unawareness, for example--but if he has no fault he cannot change for the better, but only for the worse, and for a reason which I shall discuss later, it is necessary that he must become more admirable, and not less so, at the end of the play. In other words, a hero must pass through an experience which opens his eyes to an error of his own. He must learn through suffering. In a tragedy he suffers death itself as a consequence of his fault or his attempt to correct it, but before he dies he has become a nobler person because of his recognition of his fault and the consequent alteration of his course of action. In a serious play which does not end in death he suffers a lesser punishment, but the pattern remains the same. In both forms he has a fault to begin with, he discovers that fault during the course of the action, and he does what he can to rectify it at the end. In The Green Pastures God's fault was that he believed himself perfect. He discovered that he was not perfect, and he resolved to change and grow. Hamlet's fault was that he could not make up his mind to act. He offers many excuses for his indecision until he discovers that there is no real reason for hesitation and that he has delayed out of cowardice. Lincoln, in Abe Lincoln in Illinois, has exactly the same difficulty. In the climactic scene it is revealed to him that he had hesitated to take sides through fear of the consequences to himself, and he then chooses to go ahead without regard for what may be in store for him. From the point of view of the playwright, then, the essence of a tragedy, or even of a serious play, is the spiritual awakening, or regeneration, of his hero.

When a playwright attempts to reverse the formula, when his hero makes a discovery which has an evil effect, or one which the audience interprets as evil, on his character, the play is inevitably a failure on the stage. In Troilus and Cressida Troilus discovers that Cressida is a light woman. He draws from her defection the inference that all women are faithless--that faith in woman is the possession of fools. As a consequence he turns away from life and seeks death in a cause as empty as the love he has given up, the cause of the strumpet Helen. All the glory of Shakespeare's verse cannot rescue the play for an audience, and save in Macbeth Shakespeare nowhere wrote so richly, so wisely, or with such a flow of brilliant metaphor.

For the audience will always insist that the alteration in the hero be for the better--or for what it believes to be the better. As audiences change the standards of good and evil change, though slowly and unpredictably, and the meanings of plays change with the centuries. One thing only is certain: that an audience watching a play will go along with it only when the leading character responds in the end to what it considers a higher moral impulse than moved him at the beginning of the story, though the audience will of course define morality as it pleases and in the terms of its own day. It may be that there is no absolute up or down in this world, but the race believes that there is, and will not hear of any denial.

And now at last I come to the point toward which I've been struggling so laboriously. Why does the audience come to the theatre to look on while an imaginary hero is put to an imaginary trial and comes out of it with credit to the race and to himself? It was this question that prompted my essay, and unless I've been led astray by my own predilections there is a very possible answer in the rules for playwriting which I have just cited. The theatre originated in two complementary religious ceremonies, one celebrating the animal in man and one celebrating the god. Old Greek Comedy was dedicated to the spirits of lust and riot and earth, spirits which are certainly necessary to the health and continuance of the race. Greek tragedy was dedicated to man's aspiration, to his kinship with the gods, to his unending, blind attempt to lift himself above his lusts and his pure animalism into a world where there are other values than pleasure and survival. However unaware of it we may be, our theatre has followed the Greek patterns with no change in essence, from Aristophanes and Euripides to our own day. Our more ribald musical comedies are simply our approximation of the Bacchic rites of Old Comedy. In the rest of our theatre we sometimes follow Sophocles, whose tragedy is always an exaltation of the human spirit, sometimes Euripides, whose tragi-comedy follows the same pattern of an excellence achieved through suffering. The forms of both tragedy and comedy have changed a good deal in non-essentials, but in essentials--and especially in the core of meaning which they must have for audiences--they are in the main the same religious rites which grew up around the altars of Attica long ago.

It is for this reason that when you write for the theatre you must choose between your version of a phallic revel and your vision of what mankind may or should become. Your vision may be faulty, or shallow, or sentimental, but it must conform to some aspiration in the audience, or the audience will reject it. Old Comedy, the celebration of the animal in us, still has a place in our theatre, as it had in Athens, but here, as there, that part of the theatre which celebrated man's virtue and his regeneration in hours of crisis is accepted as having the more important function. Our comedy is largely the Greek New Comedy, which grew out of Euripides' tragi-comedy, and is separated from tragedy only in that it presents a happier scene and puts its protagonist through an ordeal which is less than lethal.

And since our plays, aside from those which are basically Old Comedy, are exaltations of the human spirit, since that is what an audience expects when it comes to the theatre, the playwright gradually discovers, as he puts plays before audiences, that he must follow the ancient Aristotelian rule: he must build his plot around a scene wherein his hero discovers some mortal frailty or stupidity in himself and faces life armed with a new wisdom. He must so arrange his story that it will prove to the audience that men pass through suffering purified, that, animal though we are, despicable though we are in many ways, there is in us all some divine, incalculable fire that urges us to be better than we are.

It could be argued that what the audience demands of a hero is only conformity to race morality, to the code which seems to the spectators most likely to make for race survival. In many cases, especially in comedy, and obviously in the comedy of Moliére, this is true. But in the majority of ancient and modern plays it seems to me that what the audience wants to believe is that men have a desire to break the moulds of earth which encase them and claim a kinship with a higher morality than that which hems them in. The rebellion of Antigone, who breaks the laws of men through adherence to a higher law of affection, the rebellion of Prometheus, who breaks the law of the gods to bring fire to men, the rebellion of God in The Green Pastures against the rigid doctrine of the Old Testament, the rebellion of Tony in They Knew What They Wanted against the convention that called on him to repudiate his cuck- old child, the rebellion of Liliom against the heavenly law which asked him to betray his own integrity and make a hypocrisy of his affection, even the repudiation of the old forms and the affirmation of new by the heroes of Ibsen and Shaw, these are all instances to me of the groping of men toward an excellence dimly apprehended, seldom possible of definition. They are evidence to me that the theatre at its best is a religious affirmation, an age-old rite restating and reassuring man's belief in his own destiny and his ultimate hope. The theatre is much older than the doctrine of evolution, but its one faith, asseverated again and again for every age and every year, is a faith in evolution, in the reaching and the climb of men toward distant goals, glimpsed but never seen, perhaps never achieved, or achieved only to be passed impatiently on the way to a more distant horizon.

\newpage
\section{Whatever Hope We Have}

There is always something slightly embarrassing about the public statements of writers and artists, for they should be able to say whatever they have to say in their work, and let it go at that. Moreover, the writer or artist who brings a message of any importance to his generation will find it impossible to reduce that message to a bald statement, or even a clearly scientific statement, because the things an artist has to communicate can be said only in symbols, in the symbols of his art. The work of art is a hieroglyph, and the artist's endeavor is to set forth his vision of the world in a series of picture writings which convey meanings beyond the scope of direct statement. There is reason for believing that there is no other way of communicating new concepts save the artist's way, no other way save the artist's way of illuminating new pathways in the mind. Even the mathematician leaves the solid plane of the multiplication table and treads precariously among symbols when he advances toward ideas previously unattained.

It may be that I am trying, at this moment, to reduce to plain statement an intuitive faith of my own which cannot be justified by logic and which may lose, even for me, some of its iridescence when examined under a strong light by many searching eyes. For though the question I meant to take up was only the utility of prizes for artistic excellence, I can find no approach to that question save through a definition of the artist's faith as I see it, and no definition of that faith without an examination of the artist's place in his universe, his relation to the national culture, and the dependence of a nation on its culture for coherence and enduring significance.

Let me begin then, quite simply and honestly, even naively, with a picture of the earth as I see it. The human race, some two billion strong, finds itself embarked on a curious voyage among the stars, riding a planet which must have set out from somewhere, and must be going somewhere, but which was cut adrift so long ago that its origin is a matter of speculation and its future beyond prophecy. Our planet is of limited area, and our race is divided into rival nations and cultures that grow and press on one another, fighting for space and the products of the ground. We are ruled by men like ourselves, men of limited intelligence, with no fore-knowledge of what is to come, and hampered by the constant necessity of maintaining themselves in power by placating our immediate selfish demands. There have been men among us from time to time who had more wisdom than the majority, and who laid down precepts for the conduct of a man's brief life. Some of them claimed inspiration from beyond our earth, from spirits or forces which we cannot apprehend with our five senses. Some of them speak of gods that govern our destinies, but no one of them has had proof of his inspiration or of the existence of a god. Nevertheless there have been wise men among them, and we have taken their precepts to heart and taken their gods and their inspiration for granted.

Each man and woman among us, with a short and harried life to live, must decide for himself what attitude he will take toward the shifting patterns of government, justice, religion, business, morals, and personal conduct. We are hampered as well as helped in these decisions by every prejudice of ancestry and race, but no man's life is ready-made for him. Whether he chooses to conform or not to conform, every man's religion is his own, every man's politics is his own, every man's vice or virtue is his own, for he alone makes decisions for himself. Every other freedom in this world is restricted, but the individual mind is free according to its strength and desire. The mind has no master save the master it chooses. And each must make his choices, now as always, without sufficient knowledge and without sufficient wisdom, without certainty of our origin, without certainty of what undiscovered forces lie beyond known scientific data, without certainty of the meaning of life, if it has a meaning, and without an inkling of our racial destiny. In matters of daily and yearly living, we have a few, often fallible, rules of thumb to guide us, but on all larger questions the darkness and silence about us is complete.

Or almost complete. Complete save for an occasional prophetic voice, an occasional gleam of sci- entific light, an occasional extraordinary action which may make us doubt that we are utterly alone and completely futile in this incomprehensible journey among the constellations. From the beginning of our story, men have insisted, despite the darkness and silence about them, that they had a destiny to fulfil--that they were part of a gigantic scheme which was understood somewhere, though they themselves might never understand it. There are no proofs of this. There are only indications--in the idealism of children and young men, in the sayings of such teachers as Christ and Buddha, in the vision of the world we glimpse in the hieroglyphics of the masters of the great arts, and in the discoveries of pure science, itself an art, as it pushes away the veils of fact to reveal new powers, new laws, new mysteries, new goals for the eternal dream. The dream of the race is that it may make itself better and wiser than it is, and every great philosopher or artist who has ever appeared among us has turned his face away from what man is toward whatever seems to him most godlike that man may become. Whether the steps proposed are immediate or distant, whether he speaks in the simple parables of the New Testament or the complex musical symbols of Bach and Beethoven, the message is always to the effect that men are not essentially as they are but as they imagine and as they wish to be. The geologists and anthropologists, working hand in hand, tracing our ancestry to a humble little animal with a rudimentary fore-brain which grew with use and need, reinforce the constant faith of prophet and artist. We need more intelligence and more sensitivity if ever an animal needed anything. Without them we are caught in a trap of selfish interest, international butchery, and a creed of survival that periodically sacrifices the best to the worst, and the only way out that I can see is a race with a better brain and superior inner control. The artist's faith is simply a faith in the human race and its gradual acquisition of wisdom.

Now it is always possible that he is mistaken or deluded in what he believes about his race, but I myself accept his creed as my own. I make my spiritual code out of my limited knowledge of great music, great poetry, and great plastic and graphic arts, including with these, not above them, such wisdom as the Sermon on the Mount and the last chapter of Ecclesiastes. The test of a man's inspiration for me is not whether he spoke from a temple or the stage of a theatre, from a martyr's fire or a garden in Hampstead. The test of a message is its continuing effect on the minds of men over a period of generations. The world we live in is given meaning and dignity, is made an endurable habitation, by the great spirits who have preceded us and set down their records of nobility or torture or defeat in blazons and symbols which we can understand. I accept these not only as prophecy, but as direct motivation toward some far goal of racial aspiration. He who meditates with Plato, or finds himself shaken by Lear's ``fivefold never'? over Cordelia, or climbs the steep and tragic stairway of symphonic music, is certain to be better, both intellectually and morally, for the experience. The nobler a man's interests the better citizen he is. And if you ask me to define nobility, I can answer only by opposites, that it is not buying and selling, or betting on the races. It might be symbolized by such a figure as a farmer boy in Western Pennsylvania plowing corn through a long afternoon and saying over and over to himself certain musical passages out of Marlowe's Doctor Faustus. He might plow his corn none too well, he might be full of what we used to call original sin, but he carries in his brain a catalytic agent the presence of which fosters ripening and growth. It may be an impetus that will advance him or his sons an infinitesimal step along the interminable ascent.

The ascent, if we do climb, is so slow, so gradual, so broken, that we can see little or no evidence of it between the age of Homer and our own time. The evidence we have consists in a few mountain peaks of achievement, the age of Pericles, the centuries of Dante and Michelangelo, the reign of Elizabeth in England, the century and a half of music in Germany, peaks and highlands from which the masters seem to have looked forward into the distance far beyond our plodding progress. Between these heights lie long valleys of mediocrity and desolation, and, artistically at least, we appear to be miles beneath the upper levels traversed behind us. It must be our hope as a nation that either in pure art or in pure science we may arrive at our own peak of achievement, and earn a place in human history by making one more climb above the clouds.

The individual, the nation, and the race are all involved together in this effort. Even in our disillusioned era, when fixed stars of belief fall from our sky like a rain of meteors, we find that men cling to what central verities they can rescue or manufacture, because without a core of belief neither man nor nation has courage to go on. This is no figure of speech, no sanctimonious adjuration--it is a practical, demonstrable fact which all men realize as they add to their years. We must have a personal, a national, and a racial faith, or we are dry bones in a death valley, waiting for the word that will bring us life. Mere rationalism is mere death. Mere scientific advance without purpose is an advance toward the waterless mirage and the cosmic scavengers. The doctrine of Machiavelli is a fatal disease to the citizen or the state. The national conscience is the sum of personal conscience, the national culture the sum of personal culture--and the lack of conscience is an invitation to destruction, the lack of culture an assurance that we shall not even be remembered.

No doubt I shall be accused of talking a cloudy philosophy, of mixed metaphors and fantasy, but unless I misread my history, the artist has usually been wiser even about immediate aims than the materialist or the enthusiast for sweeping political reform. The artist is aware that man is not perfect, but that he seeks perfection. The materialist sees that men are not perfect, and erects his philosophy on their desire for selfish advantage. He fails quickly always, because men refuse to live by bread alone. The utopian sees that men seek perfection and sets out to achieve it or legislate it into existence. He fails because he cannot build an unselfish state out of selfish citizens, and he who asks the impossible gets nothing. The concepts of truth and justice are variables approaching an imaginary limit which we shall never see; nevertheless those who have lost their belief in truth and justice and no longer try for them, are traitors to the race, traitors to themselves, advocates of the dust.

To my mind a love of truth and justice is bound up in men with a belief in their destiny; and the belief in their destiny is of one piece with national and international culture. The glimpse of the godlike in man occasionally vouchsafed in a work of art or prophecy is the vital spark in a world that would otherwise stand stock still or slip backward down the grade, devoid of motive power.

For national growth and unity the artist's vision is the essential lodestone without which there is no coherence. A nation is not a nation until it has a culture which deserves and receives affection and reverence from the people themselves. Our culture in this country has been largely borrowed or sectional or local; what we need now to draw us together and make us a nation is a flowering of the national arts, a flowering of the old forms in this new soil, a renaissance of our own. If we want to live, or deserve to live, as a force or in history, we must somehow encourage the artists who appear among us, and we must encourage excellence among them. How to go about it is a problem entirely unsolved. I wish I could believe that prizes, or critics or governmental endowments were effective stimulants toward effort or excellence in any artistic field. They may be occasionally, but the greatest achievements have occurred in the absence of endowments, or professional critics or prizes, seemingly as the result of a feverish desire for accomplishment in any single art, permeating a whole society during a period long enough to allow for more than one generation of devotees. Probably an artist ean ask nothing better than a free society which likes his work and is willing to pay for it.

Looking ahead, I have no more than a hope that our nation will sometime take as great a place in the cultural history of the world as has been taken by Greece or Italy or England. So far we have, perhaps, hardly justified even the hope. But what hope there is for us lies in our nascent arts, for if we are to be remembered as more than a mass of people who lived and fought wars and died, it is for our arts that we will be remembered. The captains and the kings depart; the great fortunes wither, leaving eno trace; inherited morals dissipate as rapidly as inherited wealth; the multitudes blow away like locusts; the records and barriers go down. The rulers, too, are forgotten unless they have had the forethought to surround themselves with singers and makers, poets and artificers in things of the mind.

This is not immortality, of course. So far as I know there is no immortality. But the arts make the longest reach toward permanence, create the most enduring monuments, project the farthest, widest, deepest influence of which human prescience and effort are capable. The Greek religion is gone, but Aeschylus remains. Catholicism shrinks back toward the papal state, but the best of medieval art perishes only where its pigments were perishable. The Lutheranism of Bach retains little content for us, but his music is indispensable. And there is only one condition that makes possible a Bach, an Aeschylus, or a Michelangelo--it is a national interest in and enthusiasm for the art he practices. The supreme artist is only the apex of a pyramid; the pyramid itself must be built of artists and art-lovers, apprentices and craftsmen so deeply imbued with a love for the art they follow or practice that it has become for them a means of communication with whatever has been found highest and most admirable in the human spirit. To the young people of this country I wish to say: if you practice an art, be proud of it, and make it proud of you; if you now hesitate on the threshold of your maturity, wondering what rewards you should seek, wondering perhaps whether there are any re- wards beyond the opportunity to feed and sleep and breed, turn to the art which has moved you most readily, take what part in it you can, as participant, spectator, secret practitioner, or hanger-on and waiter at the door. Make your living any way you can, but neglect no sacrifice at your chosen altar. It may break your heart, it may drive you half mad, it may betray you into unrealizable ambitions or blind you to mercantile opportunities with its wandering fires. But it will fill your heart before it breaks it; it will make you a person in your own right; it will open the temple doors to you and enable you to walk with those who have come nearest among men to what men may sometime be. If the time arrives when our young men and women lose their extravagant faith in the dollar and turn to the arts we may then become a great nation, nurturing great artists of our own, proud of our own culture and unified by that culture into a civilization worthy of our unique place on this rich and lucky continent between its protecting seas.

\newpage
\section{A Prelude to Poetry in the Theatre}

Experiment in the theatre is made difficult and expensive by the fact that a play must find an audience at once or have no chance of finding one later. There is no instance in the theatre of a writer who left behind him a body of unappreciated work which slowly found its public, as, for example, the work of Shelley and Keats found a belated public after they had left the scene. It follows that the playwright must pluck from the air about him a fable which will be of immediate interest to his time and hour, and relate it in a fashion acceptable to his neighbors. That is the job for which he is paid. But he will also try to make that fable coincide with something in himself that he wants to put in words. A certain cleverness in striking a compromise between the world about him and the world within has characterized the work of the greatest as well as the least of successful playwrights, for they must all take an audience with them if they are to continue to function. Some may consider it blasphemy to state that this compromise must be a considered and conscious act-- will believe that the writer should look in his heart and write--but in the theatre such an attitude leaves the achievement entirely to chance, and a purely chance achievement is not an artistic one.

Yet when a writer sits down before white paper to make this necessary compromise he finds himself alone among imponderables. Nobody has ever known definitely what any audience wanted. A choice must be made with only intuition and a mass of usually irrelevant information as the guides. One who thinks more of his job than his fame will therefore play safe by repressing his personal preferences and going all the way in the direction of what he believes the public wants. One who thinks as much of his fame as of his job will often hope the public is ready for a theme only because he wishes to treat it--or ready for a dramatic method only because he wishes to employ it. I may have been somewhat guilty of this last misapprehension in Winterset, for I have a strong and chronic hope that the theatre of this country will outgrow the phase of journalistic social comment and reach occasionally into the upper air of poetic tragedy. I believe with Goethe that dramatic poetry is man's greatest achievement on his earth so far, and I believe with the early Bernard Shaw that the theatre is essentially a cathedral of the spirit, devoted to the exaltation of men, and boasting an apostolic succession of inspired high priests which extends further into the past than the Christian line founded by St. Peter. It has been, even at its best, a democratic temple, decorated with more gargoyles than saints, generously open to wits, clowns, excoriating satirists, false prophets and crowds of money-changers with a heavy investment in the mysteries. Lately it has recognized the mysteries only as a side-show, and has been over-run with guides who prove to an eager public that all saints are plaster and all prophets fakes.

When Shaw began his furious critical assault on the romantic theatre which was lingering out the last decade of the nineteenth century he became the prophet of the theatre of realistic social protest which he has dominated in England during his life-time and which is still the controlling pattern for plays of the English-speaking stage. A few original and outstanding playwrights, J. M. Synge, Sean O'Casey and Eugene O'Neill among them, may seem to fall completely outside the Shavian category. Synge was too poetic, symbolic and savage to work in any social harness, O'Casey's world-searing irony is too tremendous to come under the head of protest, and O'Neill has been seeking an escape from realism throughout his whole career. All three went beyond Shaw's amusing balance and clarity into a tragic, sinister and often brutal world. Synge wrote of that world with extraordinary beauty and a lethal precision; O'Casey and O'Neill write of it with passion and a sometimes choking defiance. O'Casey, in addition, has a power over words which lifts such a play as The Plough and the Stars out of its local setting and makes it, in my opinion, the finest contemporary play I have seen in the theatre.

But none of these three, nor Shaw himself, has written plays which we can set unquestioningly beside the best we can pick up in the library--and the reason for that is a fairly simple one. Our modern dramatists (with the exception of O'Casey, who has lately tried his 'prentice wings) are not poets, and the best prose in the world is inferior on the stage to the best poetry. It is the fashion, I know, to say that poetry is a matter of content and emotion, not of form, but this is said in an age of prose by prose writers who have not studied the effect of form on content or who wish to believe there is no limit to the scope of the form they have mastered. To me it is inescapable that prose is the language of information and poetry the language of emotion. Prose can be stretched to carry emotion, and in some exceptional cases, as in Synge's and O'Casey's plays, can occasionally rise to poetic heights by substituting the unfamiliar speech rhythms of an untutored people for the rhythm of verse. But under the strain of an emotion the ordinary prose of our stage breaks down into inarticulateness, just as it does in life. Hence the cult of understatement, hence the realistic drama in which the climax is reached in an eloquent gesture or a moment of meaningful silence.

The majority of present-day play-goers have never seen any other kind of play and see no reason why there should be any other. So emphatic is this feeling that one is doubtful of being able to explain to this majority that verse was once the accepted convention on the stage, as prose is now, that prose fought its way into the play-books with difficulty at the beginning of the scientific era in which we live and will hold its place there only so long as men make a religion of fact and believe that information, conveyed in statistical language, can make them free.

For the stage is still a cathedral, but just now a journalistic one, dominated by those who wish to offer something immediate about our political, social or economic life. Like every other existing condition it gives the illusion of permanence, but it will change. An age of reason will be followed once more by an age of faith in things unseen. The cathedral will then house the mysteries again, along with the jugglers and the vendors of rose-colored spectacles. What faith men will then have, when they have lost their certainty of salvation through laboratory work, I don't know, having myself only a faith that men will have a faith. But that it will involve a desire for poetry after our starvation diet of prose I have no doubt. Men have not been altered by the invention of airplanes and the radio. They are still alone and frightened, holding their chance tenure of life in utter isolation in this desolate region of revolving fires. Science may answer a few necessary questions for them, but in the end science itself is obliged to say that the fact is created by the spirit, not spirit by the fact. Our leading scientists are already coming to this conclusion, rather reluctantly and with some surprise.

Unless I am greatly mistaken many members of the theatre audience have anticipated this conclusion by one of those intuitional short-cuts which confound ``the devotees of pure reason, and are not only ready but impatient for plays which will take up again the consideration of man's place and destiny in prophetic rather than prosaic terms. It is incumbent on the dramatist to be a poet, and incumbent on the poet to be prophet, dreamer and interpreter of the racial dream. Men have come a long way from the salt water in the millions of years that lie behind them, and have a long way to go in the millions of years that lie ahead. We shall not always be as we are-- but what we are to become depends on what we dream and desire. The theatre, more than any other art, has the power to weld and determine what the race dreams into what the race will become. All this may sound rather far-fetched in the face of our present Broadway, and Broadway may laugh at it unconscionably, but Broadway is itself as transient as the real-estate values under its feet. Those of us who fail to outlive the street in which we work will fail because we have accepted its valuations and measured our product by them. For though on the surface we are still a pioneer people, ashamed of aspiration, offended by the deliberate quest for beauty, able to accept beauty only when it seems achieved by accident, our pioneer days are over and we must set about moulding ourselves at least one art form worthy of the leading nation of the world or be set down finally as barbarians and carry that name with us into the darkness to which all nations sooner or later descend. Our theatre is the one really living American art. It has size, vitality and popular interest. But it is still in the awkward and self-conscious age, concealing its dreams by clowning, burlesquing the things it most admires. Those who have read their literary history carefully know that now is the time for our native amusements to be transformed into a national art of power and beauty. It needs the touch of a great poet to make the transformation, a poet comparable to Aeschylus in Greece or Marlowe in England. Without at least one such we shall never have a great theatre in this country, and he must come soon, for these chances don't endure forever.

I must add, lest I be misunderstood, that I have not mistaken myself for this impending phenomenon. I have made my living as teacher, journalist and playwright and have only that skill as a poet which may come from long practice of an art I have loved and studied and cannot let alone. When I wrote my first play, White Desert, I wrote it in verse because I was weary of plays in prose that never lifted from the ground. It failed, and I did not come back to verse again until I had discovered that poetic tragedy had never been successfully written about its own place and time. There is not one tragedy by Aeschylus, Sophocles, Euripides, Shakespeare, Corneille or Racine which did not have the advantage of a setting either far away or long ago. With this admonition in mind I wrote Elizabeth the Queen and a succession of historical plays in verse, some of them successful, and found myself immediately labelled a historical and romantic playwright, two terms I found equally distasteful. Winterset is largely in verse, and treats a contemporary tragic theme, which makes it more of an experiment than I could wish, for the great masters themselves never tried to make tragic poetry out of the stuff of their own times. To do so is to attempt to establish a new convention, one that may prove impossible of acceptance, but to which I was driven by the lively historical sense of our day--a knowledge of period, costume and manners which almost shuts off the writer on historical themes from contemporary comment. Whether or not I have solved the problem in Winterset is probably of little moment. But it must be solved if we are to have a great theatre in America. Our theatre has not yet produced anything worthy to endure--and endurance, though it may be a fallible test, is the only test of excellence.

\newpage
\section{The Politics of Knickerbocker Holiday}

Knickerbocker holiday was obviously written to make an occasion for Kurt Weill's music, and since Mr. Weill responded by writing the best score in the history of our theatre, and since the public has voted an emphatic approval at the box office, the whole venture would seem to justify itself without further comment.

But there has been a good deal of critical bewilderment over the political opinions expressed in the play, and not a little resentment at my definitions of government and democracy. I should like to explain that it was not my intention to say anything new or shocking on either subject, but only to remind the audience of the attitude toward government which was current in this country at the time of the revolution of 1776 and throughout the early years of the Republic. At that time it was generally believed, as I believe now, that the gravest and most constant danger to a man's life, liberty and happiness is the government under which he lives. It was believed then, as I believe now, that a civilization is a balance of selfish interests, and that a government is necessary as an arbiter among these interests, but that the government must never be trusted, must be constantly watched, and must be drastically limited in its scope, because it, too, is a selfish interest and will automatically become a mo- nopoly in crime and devour the civilization over which it presides unless there are definite and positive checks on its activities. The Constitution is a monument to our forefathers' distrust of the state, and the division of powers among the legislative, judicial and executive branches has succeeded so well for more than a century in keeping the sovereign authority in its place that our government is now widely regarded as a naturally wise and benevolent institution, capable of assuming the whole burden of social and economic justice.

The thinking behind our Constitution was dominated by such men as Franklin and Jefferson, men with a high regard for the rights of the individual, combined with a cold and realistic attitude toward the blessings of central authority. Knowing that government was a selfish interest, they treated it as such, and asked of it no more than a selfish interest can give. But the coddled young reformer of our day, looking out on his world, finding merit often unrewarded and chicanery triumphant, throws prudence to the winds and grasps blindly at any weapon which seems to him likely to destroy the purse-proud haves and scatter their belongings among the deserving have-nots. Now he is right in believing that the accumulation of too much wealth and power in a few hands is a danger to his civilization and his liberty. But when the weapon he finds is a law, and when the law he enacts increases the power of the government over men's destinies, he is fighting a les- ser tyranny by accepting a greater and more deadly one, and he should be aware of that fact.

A government is always, as Stuyvesant says in Knickerbocker Holiday, ``a group of men organized to sell protection to the inhabitants of a limited area at monopolistic prices.'' The members of a government are not only in business, but in a business which is in continual danger of lapsing into pure gangsterism, pure terrorism and plundering, buttered over at the top by a hypocritical pretense at patriotic unselfishness. The continent of Europe has been captured by such governments within the last few years, and our own government is rapidly assuming economic and social responsibilities which take us in the same direction. Whatever the motives behind a government-dominated economy, it can have but one result, a loss of individual liberty in thought, speech and action. A guaranteed life is not free. Social security is a step toward the abrogation of the individual and his absorption into that robot which he has invented to serve him--the paternal state.

When I have said this to some of the youthful proponents of guaranteed existence I have been met with the argument that men must live, and that when the economic machinery breaks down men must be cared for lest they starve or revolt. This is quite true and nobody opposes emergency relief. It is the attempt to make the emergency and the relief permanent that constitutes an attack on our free institutions. The greatest enemies of democracy, the most violent reactionaries, are those who have lost faith in the capacity of a free people to manage their own affairs and wish to set up the government as a political and social guardian, running their business and making their decisions for them.

For life is infinitely less important than freedom. A free man has a value to himself and perhaps to his time; a ward of the state is useless to himself--useful only as so many foot-pounds of energy serving those who manage to set themselves above him. A people which has lost its freedom might better be dead, for it has no importance in the scheme of things except as an evil power behind the word of a dictator. In our hearts we all despise the man who wishes the state to take care of him, who would not rather live meagerly as he pleases than suffer a fat and regimented existence. Those who are not willing to sacrifice their lives for their liberty have never been worth saving. Throughout remembered time every self-respecting man has been willing to defend his liberty with his life. If our country goes totalitarian out of a soft-headed humanitarian impulse to make life easy for the many, we shall get what we vote for and what we deserve, for the choice is still before us, but we shall have betrayed the race of men, and among them the very have-nots whom we subsidize. Our western continent still has the opportunity to resist the government-led rush of barbarism which is taking Europe back toward Attila, but we can only do it by running our government, and by refusing to let it run us.

If the millions of workingmen in this country who are patiently paying their social security dues could glimpse the bureaucratic absolutism which that act presages for themselves and their children they would repudiate the whole monstrous and dishonest business overnight. When a government takes over a people's economic life it becomes absolute, and when it has become absolute it destroys the arts, the minds, the liberties and the meaning of the people it governs. It is not an accident that Germany, the first paternalistic state in Europe, should be governed now by an uncontrollable dictator; not an accident that Russia, adopting a centrally administered economy for humanitarian reasons, should arrive at a tyranny bloodier and more absolute than that of the Czars. Men who are fed by their government will soon be driven down to the status of slaves or cattle.

All this was known to the political leaders who put our Constitution together after the revolution against England. The Constitution is so built that while we adhere to it we cannot be governed by one man or one faction, and when we have made mistakes we reserve the right to change our minds. The division of powers and the rotation of offices was designed to protect us against dictatorship and arbitrary authority. The fact that there are three branches of government makes for a salutary delay and a blessed inefficiency, the elective rotation makes for a government not by cynical professionals, but by normally honest and fairly incompetent amateurs. That was exactly what the wary old founding fathers wanted, and if we are wise we shall keep it, for no scheme in the history of the world has succeeded so well in maintaining the delicate balance between personal liberty and the minimum of authority which is necessary for the free growth of ideas in a tolerant society.

\newpage
\section{Yes, by the Eternal}

In his estimate of a season's poetic plays not so long ago Mr. Max Eastman over-praised my playwriting most egregiously--and naturally I tried to forgive him. He also damned some of my verse writing for the theatre most engagingly--and to that I had no answer. Nobody but an egomaniac can defend his own verse. But when Mr. Eastman stepped off the deep end and began to discuss the poetic drama in general he revealed a confusion of mind so definite and so nearly typical of modern critics of the theatre that I am tempted to point it out.

During the last quarter century, as he says, Mr. Eastman has been prophesying and hoping for an imaginative rebirth in this country which would include a re-establishment of the poetic theatre. But the poetic theatre he wants is a special kind of poetic theatre, one which will throw its influence behind ``the intelligent, which is to say the practical scientific--which is to say the real--effort to solve the problems of life on this planet.'' The italics are Mr. Eastman's. And the italics leave no doubt that Mr. Eastman is asking not only for a new kind of poetic theatre, but a new kind of poetry as well. Never in the history of the world has poetry of any excellence thrown its weight toward the practical or scientific reorganization of the affairs of men. Poetry is just as unfit for that business as for making up the accounts of a brokerage house. As for plays, even a play in prose loses its franchise over an audience the moment it begins to discuss the blueprints for an almost perfect state. Satire in verse and satire in prose are different matters. An attack on the existing order sometimes animates good poetry or a good play. But satire at its best is second or third best, and obviously Mr. Eastman was not asking for satire. He was not even asking for the burning anger at injustice which sometimes fires a poem or play. He was asking for something immediately constructive, scientifically and practically constructive.

He will never get it in poetry either on or off the stage, because a concrete constructive policy in the affairs of our planet has never yet called forth great poetry, and by its very nature, cannot invoke great poetry, or any poetry at all.``Poetry is,'' as Mr. Eastman says, ``a way of using language.'' But it is a way of using language that impels the user powerfully toward emotional utterance, impels him away from the small change of political economy and toward whatever vision he may be able to formulate of human destiny. The poetic impulse is a mystery that Mr. Eastman cannot explain any more than I can, but we can be certain of this about it, that it is no timeserver, that it cannot be hitched to the Marxian plow or turned on to campaign for a single-tax. It is indulged in by dreamers, and not by practical men, yet wherever the practical men have gone in search of truth or wealth they have found a few errant dream- ers ahead of them. The poet is usually long dead by the time he inherits the earth, but he does inherit it.

The great poetry of Greece, of Italy and of England is nearly all as mystic in concept and as prophetic in tone as the Old Testament itself. Prophetic with the eye on the distant horizon, not on the excavation in the foreground. Over and over again the poet, over-occupied with the horizon, steps, still singing, into the excavation for a new edifice, and is walled into the foundation quite unheeded. Over and over again his musical cry is remembered when the building is ruined or effaced; often there is little known of an entire civilization save the words of an obscure singer long ago buried under the fallen walls of a lost and forgotten political order.

And if we examine these musical cries for their meaning we find that the writers of epic, of lyric and dramatic verse are alike in having no pressing plans for the race. The writers of epics celebrate the youth, the hope, the victory, the disillusion and the defeat of men; the writers of lyrics are always young, and their constant theme is the anguish of youth in its first contact with reality and inevitable despair; the authors of tragedy offer the largest hope for mankind which I can discern in the great poetry of the earth, a hope that man is greater than his clay, that the spirit of man may rise superior to physical defeat and death. The theme of tragedy has always been victory in defeat, a man's conquest of himself in the face of annihilation. The last act of a tragedy contains the moment when the wheel of a man's fate carries him simultaneously to spiritual realization and to the end of his life. The message of tragedy is simply that men are better than they think they are, and this message needs to be said over and over again in every tongue lest the race lose faith in itself entirely. Perhaps Mr. Eastman thinks this message old-fashioned. Perhaps he still wants the Muses to divest themselves of whatever misty stuff they are wearing, don overalls, and go to work in the factory of the mind. But he has not said what practical and scientific plans for the advancement of the race he would set them to work on. I know of no such plans which would not be far more accurately discussed in cold prose in print than in verse on the stage. In fact, I know of no such plans which look further ahead for the race than the next few years, and none which appears so convincing in project or operation that a poet might be tempted to trail clouds of glory over it. Knowing that Mr. Eastman is a Socialist I may be pardoned for suspecting that what he really wants is a poetic theatre devoted to Marxism, and being aware that this is a contradiction in terms I can only ask his pardon for suspecting him of demanding the obviously impossible. (Perhaps I may be pardoned also for remarking that, though Mr. Eastman is highly scornful of any ``faith in things unseen,'' surely Marxism is such a faith, for the Marxist state has never been seen, and can only be hoped for by its devotees.)

At this point Mr. Eastman, if he has been so far patient, will again call me a defeatist, for defeatist is the label which the Marxists have devised for all who refuse to fight with their weapons. But I am not a defeatist. My hope for the human race is that it will so far improve in mentality and magnanimity, over a period of millenniums, that it will be able to govern itself without recourse to violence. It seems to me that it requires more steady courage to hold such a belief than to hide under the childish illusion that mankind will save itself by clever inventions, political, mechanical, or scientific. Marxism is an ingenious political invention which would work if the majority of men were honest and unselfish, but they are not. They may be later on. My hope is that they will be. What savage disciplines, what devotion to the arts and sciences, what birth and death of races and race moralities must intervene between us and that far consummation, no man can guess or estimate. But when the race gets there, if it does, the poets will still be ahead of it, examining a still more distant future. For what the poets are always asking for, visioning, and projecting is man as he must and will be, man a step above and beyond his present, man as he may be glimpsed on some horizon of dream, a little nearer what he himself wishes to become.

\end{document}
